\documentclass[10pt, aspectratio=169]{beamer}
\usepackage{../beamer_theme_408}

\title{408考研零基础 --- 数据结构}
\subtitle{1-5 链表进阶}
\author{}
\date{}


\begin{document}


\begin{frame}{本节目标}
    \begin{enumerate}
        \item 理解循环链表的结构特征及其与单链表的区别
        \item 理解双向链表的结点结构与插入删除操作
        \item 能够根据实际需求选择顺序表或链表
    \end{enumerate}
\end{frame}

\begin{frame}{循环链表}
    \textbf{特征}：最后一个结点的指针域指向头结点，形成环
    \vspace{0.5em}
    \begin{itemize}
        \item 从任意结点出发可访问所有结点
        \item 判空条件：$\text{head} \to \text{next} == \text{head}$
        \item 遍历终止条件：$p == \text{head}$
    \end{itemize}
    \vspace{0.5em}
    \textbf{尾指针优化}
    \begin{itemize}
        \item 使用尾指针 $\text{rear}$ 标识循环链表
        \item 头结点：$\text{rear} \to \text{next}$
        \item 尾结点：$\text{rear}$ 本身
        \item 可在 $O(1)$ 时间内访问头尾
    \end{itemize}
\end{frame}

\begin{frame}[fragile]{双向链表的结构}
    \textbf{结点包含三个部分}
    \begin{itemize}
        \item 前驱指针 prior
        \item 数据域 data
        \item 后继指针 next
    \end{itemize}
    \vspace{0.5em}
    \begin{lstlisting}[language=C, basicstyle=\ttfamily\small]
typedef struct DNode {
    ElemType data;
    struct DNode *prior, *next;
} DNode, *DLinkList;
    \end{lstlisting}
    \vspace{0.5em}
    \begin{block}{优势}
        支持双向遍历，查找前驱 $O(1)$
    \end{block}
\end{frame}

\begin{frame}{双向链表 --- 插入操作}
    \textbf{在结点 $p$ 之后插入结点 $s$}
    \vspace{0.5em}
    \begin{enumerate}
        \item $s \to \text{next} = p \to \text{next}$
        \item $p \to \text{next} \to \text{prior} = s$（注意判空）
        \item $s \to \text{prior} = p$
        \item $p \to \text{next} = s$
    \end{enumerate}
    \vspace{0.5em}
    \begin{alertblock}{注意}
        若 $p$ 是最后一个结点，$p \to \text{next}$ 为 NULL，步骤 2 跳过
    \end{alertblock}
\end{frame}

\begin{frame}{双向链表 --- 删除操作}
    \textbf{删除 $p$ 的后继结点 $q$}
    \vspace{0.5em}
    \begin{enumerate}
        \item $p \to \text{next} = q \to \text{next}$
        \item $\text{if}(q \to \text{next} \ne \text{NULL}) \; q \to \text{next} \to \text{prior} = p$
        \item $\text{free}(q)$
    \end{enumerate}
    \vspace{0.5em}
    时间复杂度 $O(1)$
\end{frame}

\begin{frame}{顺序表与链表对比}
    \begin{center}
    \begin{tabular}{|c|c|c|}
        \hline
        \textbf{维度} & \textbf{顺序表} & \textbf{链表} \\
        \hline
        存储方式 & 连续 & 离散 \\
        \hline
        随机访问 & $O(1)$ & $O(n)$ \\
        \hline
        插入删除 & $O(n)$ & $O(1)$ \\
        \hline
        存储密度 & 1 & $< 1$ \\
        \hline
        空间分配 & 预分配 & 动态分配 \\
        \hline
    \end{tabular}
    \end{center}
\end{frame}

\begin{frame}{本节总结}
    \begin{enumerate}
        \item 循环链表：尾指针指向头结点，形成环
        \item 双向链表：prior + next 双向指针
        \item 插入删除需修改四个指针，注意边界
        \item 顺序表适合查找密集、规模可预估场景
        \item 链表适合插入删除密集、规模不确定场景
    \end{enumerate}
\end{frame}

\end{document}
